\subsection{Positive mass with a degenerate smooth residual}
\label{sec:degenerate-projective-residuals}
\index{Reference stability!degenerate residual}

A second family has a logarithmic pole and a smooth resolved residual that degenerates along a different exceptional divisor. Fix $p=[1:0:\cdots:0]$ and $q\ne p$ on $\mathbb P^n$. On $X=\operatorname{Bl}_q\mathbb P^n$ keep $\Theta=\pi^*\omega_{\mathrm{FS}}$. For $0<t<1$, put
\[
 \psi_t=t\log\frac{\sum_{i=1}^n|Z_i|^2}{\sum_{i=0}^n|Z_i|^2},
 \qquad K_t=\{p\in\Sigma:\textstyle\sum_i p_i\geq t\}.
\]
The local current potential of $\psi_t$ is $(1-t)L+t\log\sum_i\mathrm{e}^{x_i}$. Its bounded-distance support function is
\[
 (1-t)h_\Sigma+t\max_i x_i=h_{K_t},\qquad
 K_t=(1-t)\Sigma+t\operatorname{conv}(e_1,\ldots,e_n).
\]
Thus $\psi_t\sim\psi_{K_t}\sim\Phi_{K_t}$, and $n!\operatorname{Vol}_n(K_t)=1-t^n$.

\begin{proposition}\label{prop:degenerate-projective-residual}
Let $\tau:Z\to X$ resolve the pole by blowing up the point above $p$. There is a smooth semipositive residual $\alpha_t$ such that
\begin{equation}\label{eq:degenerate-projective-residual}
 (\pi\circ\tau)^*\omega_{\mathrm{FS}}
       +\ddc(\psi_t\circ\pi\circ\tau)=\alpha_t+t[E_p].
\end{equation}
It is positive away from $E_q$, has rank exactly one at every point of $E_q$, and has top integral $1-t^n>0$. Every ordinary-type solution associated with $\pi^*\Phi_{K_t}$ lifts, after subtracting the prescribed divisor pole, to a bounded canonical $\alpha_t$-plurisubharmonic remainder.
\end{proposition}
\begin{proof}
On $Y=\operatorname{Bl}_p\mathbb P^n$, let $\pi_p$ be the blowdown and $Q_p:Y\to\mathbb P^{n-1}$ the resolved direction map. The incidence graph gives the K\"ahler form
\[
 \beta_t=(1-t)\pi_p^*\omega_{\mathrm{FS}}
             +tQ_p^*\omega_{\mathrm{FS},n-1}.
\]
The disjoint centers identify $Z$ with the blowup $\sigma:Z\to Y$ at the point above $q$. Put $\alpha_t=\sigma^*\beta_t$. Poincar\'e--Lelong gives \eqref{eq:degenerate-projective-residual}. There is no $E_q$ coefficient, since $\psi_t$ is locally bounded at $q$.

In a blowup chart $\sigma$ is $(s,w_2,\ldots,w_n)\mapsto(s,sw_2,\ldots,sw_n)$. At $s=0$ its differential has rank one and kills the $n-1$ tangent directions of $E_q$. Positivity of $\beta_t$ gives the exact rank assertion; off $E_q$ the map is biholomorphic.

Write $H=\pi_p^*[\omega_{\mathrm{FS}}]$ and $E=[E_p]$. The current identity gives $[\beta_t]=H-tE$. Since $H|_{E_p}=0$, all mixed top intersections vanish. The normal line of $E_p$ is $\mathcal O_{\mathbb P^{n-1}}(-1)$, so $\int E^n=(-1)^{n-1}$. Therefore
\[
 \int_Z\alpha_t^n=\int_Y(H-tE)^n=1-t^n.
\]
The first equality uses the degree-one blowup. An ordinary-type solution differs boundedly from the pulled-back prototype off $E_p$. After adding a local smooth residual potential, its bounded remainder is plurisubharmonic there and extends canonically across $E_p$ by \cref{lem:qpsh-removable-analytic}. This proves the last assertion without inserting a divisor value into an unbounded product.
\end{proof}

\begin{corollary}\label{cor:degenerate-residual-reference-stability}
Suppose $t_j\to t$ in $[t_{\min},t_{\max}]\subset(0,1)$, and let $r_j>0$ tend to $r>0$. The models $r_j\pi^*\Phi_{K_{t_j}}$ satisfy \cref{thm:convex-body-reference-stability}, including ordinary-type solvability at the limit, and
\[
 \sup_G|B_j(G)-B(G)|
 \leq n|\log(r_j/r)|+
       \frac{nt_{\max}^{n-1}}{1-t_{\max}^n}|t_j-t|.
\]
All continuous-test supports are the normalized exponential probabilities from that theorem. The zero-mass endpoint $t=1$ admits no positive-density equation in the prototype's ordinary type.
\end{corollary}
\begin{proof}
Choose $p_{*,i}=t_{\max}/n$ and $\delta=(1-t_{\max})/2$. Their translated simplex lies in every $K_t$ in the interval. For $s\geq t$, move a point of $K_t$ toward a vertex until its coordinate sum reaches $s$. This gives
\[
 d_H(K_s,K_t)\leq\frac{\sqrt2}{1-t_{\max}}|s-t|.
\]
Apply the stability theorem and use the derivative bound for $\log(1-t^n)$ to obtain the sharper displayed functional estimate.

At $t=1$, the resolved residual is pulled back from $\mathbb P^{n-1}$, so its top integral is zero. Every ordinary-type potential has a bounded residual-plurisubharmonic remainder. Bounded closed-product Stokes gives the same zero top mass for that remainder; positivity forces its top measure to vanish. Thus it cannot solve $c\mathrm{e}^F\mathrm{d}V$ with bounded $F$ and $c>0$. In contrast, the endpoint $a=1$ in \cref{cor:curved-reference-stability} has positive mass and remains solvable.
\end{proof}

Scaling by $r>0$ multiplies the resolved residual and the prescribed divisor coefficient by $r$. Solvability is established on $X$ by the model criterion, without treating the resolved residual as strictly positive. These criteria use closed projective backgrounds and fixed modifications. Convergence for general nonclosed Hermitian backgrounds, moving modifications or unrestricted singularity degeneration remains open. The supporting probabilities here identify the limiting functional; no convergence of solution potentials or rough weighted adjoint products is inferred.
